\section*{الفصل \ref{ch:b3:covariant-electromagnetism} --- الكهرطيسية التغايرية}

\begin{solution}{exo:b3:covariant-electromagnetism:1}
(أ) $\gamma = 1.25$: ومنه $a'^0 = 1.25(5 - 1.8) = 4.0$ و $a'^1 = 1.25(3 -
3) = 0$؛ و $a\cdot a = 25 - 9 = 16 = 16 - 0$. (ب) التحويل
خطي، فيتوزع على المجاميع. (ج) $(c, \vect v)$: لا --- فمركبتها
``الزمنية'' واحدة لكل الجسيمات في حين أن الاختلاط
يقتضي غير ذلك؛ أما $\gamma(c, \vect v) = \dd x^\mu/\dd\tau$ فنعم، إذ هي
رباعية مقسومة على المقدار الثابت $\dd\tau$. (د) مركبات $\vect E$ الثلاث
تختلط بمركبات $\vect B$، لا بأيّ سلّمي: فالمجالان يملآن
موتّرًا ذا دليلين، لا رباعية.
\end{solution}

\begin{solution}{exo:b3:covariant-electromagnetism:2}
(أ) $u = I/nSe = 10/(\num{8.5e28} \times \num{e-6} \times
\num{1.6e-19}) = \qty{7.4e-4}{m/s}$. (ب) الشبكة: $(nec, \vect 0)$؛
والإلكترونات: $(-nec, -ne\vect u)$ --- فكثافة تيارها
$\rho_-\vect u = -ne\vect u$ تشير عكس انجرافها، وذاك
\emph{هو} اتجاه التيار الاصطلاحي. (ج) كلاهما ساكن ومنتظم: فينعدم كل حدّ. (د)
المجموع: $(0, \vect\jmath)$ مع $j = I/S = \qty{e7}{A/m^2}$ --- أي
تيار صرف بلا شحنة، وهو التشكيل الذي يجعل نسبية السلك
دقيقة الالتباس.
\end{solution}

\begin{solution}{exo:b3:covariant-electromagnetism:3}
(أ) $E_y$: من $F^{20} = E_y/c$؛ و $B_x$: من $F^{32} = B_x$. (ب) من أجل $\vect B
= B\vect e_z$، لا يبقى غير $F^{12} = -B$ و $F^{21} = +B$؛ ومن أجل الموجة،
$F^{20} = E/c$ و $F^{12} = -E/c$ (مع قرينيهما المتخالفَي التناظر). (ج)
$F^{12} = \partial^1A^2 - \partial^2A^1 = -\partial_xA_y +
\partial_yA_x = -(\operatorname{\vect{curl}}\vect A)_z = -B_z$:
وهو يطابق المصفوفة. (د) معدل تغير طاقة الجسيم
هو $\propto F^{0\nu}u_\nu$؛ وفي مرجع السكون $u_\nu = (c, \vect 0)$ ويجعل
تخالف التناظر $F^{00} = 0$: أي قوة لا يمكن أن تؤدي شغلًا على
جسيم ساكن --- وهي الخاصية التي تعرّف الجزء المغناطيسي.
\end{solution}

\begin{solution}{exo:b3:covariant-electromagnetism:4}
(أ) تعزيز على منحى المجال: فتبقى المركبات الطولية على حالها، $\vect
E' = E_0\vect e_y$ و $\vect B' = \vect 0$. (ب) وتعزيز على المنحى $\vect
e_x$: فيكون $E'_y = \gamma E_0$ و $B'_z = -\gamma vE_0/c^2$ --- إذ تنساب الصفيحتان
الآن تيارين سطحيين $\pm\sigma'v$، وصفيحتا تيار متعاكستان
تحصران مجالًا كهذا بالضبط. (ج) $E'^2 - c^2B'^2 =
\gamma^2E_0^2(1 - \beta^2) = E_0^2$؛ و $\vect E'\cdot\vect B' = 0$:
فكلاهما محفوظ. (د) تتقلص الصفيحتان على المحور $x$، ومنه $\sigma' =
\gamma\sigma$، وهو متسق مع $E' = \sigma'/\varepsilon_0 = \gamma
E_0$؛ أما التباعد، وهو مستعرض على التعزيز، فلا يتغير.
\end{solution}

\begin{solution}{exo:b3:covariant-electromagnetism:5}
(أ) تدفع القوة المغناطيسية $-e\vect v\wedge\vect B$ الإلكترونات
على طول القضيب: فالقوة المحركة $= vBL$. (ب) وفي مرجع القضيب يكون القضيب
ساكنًا --- فلا قوة مغناطيسية على شحنات ساكنة --- لكن
التحويل يقدّم $\vect E' = \gamma\,\vect v\wedge\vect B$:
أي مجال كهربائي صريح يقوم بالسَّوق. (ج) $E'L = \gamma vBL
\approx vBL$ في المرتبة $v/c$. (د) كان لا بدّ أن ينجح لأن القوة المحركة
``الحركية'' و``المحوّلية'' ظاهرة واحدة مقروءة في مرجعين:
فقاعدة التدفق هي التغايرية بثياب 1831 --- ويفتتح بحث أينشتاين
سنة 1905 بهذا اللاتناظر بين المغنط والناقل بالضبط.
\end{solution}

\begin{solution}{exo:b3:covariant-electromagnetism:6}
(أ) المركبات على المحور $x$ لا تُمسّ؛ وأما المستعرضة فلدينا
$E_y'^2 + E_z'^2 - c^2(B_y'^2 + B_z'^2) = \gamma^2[(E_y - vB_z)^2 +
(E_z + vB_y)^2 - c^2(B_y + vE_z/c^2)^2 - c^2(B_z - vE_y/c^2)^2]$؛
فتتقاصّ الحدود المتقاطعة وتجمع المربعات $\gamma^2(1 -
\beta^2) = 1$ مضروبًا في التركيب غير المحوَّل. (ب) والمحاسبة
نفسها على $E_xB_x + E_yB_y + E_zB_z$. (ج) $E^2 - c^2B^2 =
3c^2B^2 > 0$: فنعزّز على منحى $\vect E\wedge\vect B$ بالسرعة $v = c^2B/E =
c/2$؛ فيكون هناك $B' = 0$ والمجال الباقي $E' = E/\gamma =
\sqrt3\,cB$ --- ومربعه يعطي الثابت، كما يجب. (د) للموجة الضوئية ثابتان
معدومان: فأيّ مرجع لا بدّ أن يعيد $E' = cB' \perp$، ولا يعطي مجالًا صرفًا أبدًا.
\end{solution}

\begin{solution}{exo:b3:covariant-electromagnetism:7}
(أ) مع $v_{\text{d}} = E/B$: يكون $E'_y = \gamma(E - v_{\text{d}}B) =
0$؛ و $B'_z = \gamma(B - v_{\text{d}}E/c^2) = B/\gamma$: أي مجال
مغناطيسي صرف موهَن قليلًا. (ب) هناك: دائرة عند $qB'/\gamma m$؛ وعودةً
إلى المخبر: تلك الدائرة مضافًا إليها الانجراف المنتظم --- أي دويري يزحف
بالسرعة $E/B$ عموديًا على المجالين، وهو مسار الشحنات في
المغنطرون. (ج) الجسيم المتحرك بالسرعة $v_{\text{d}}$ بالضبط ساكن
في مرجع الانجراف، حيث المجال الوحيد مغناطيسي ولا
يحسّ شيئًا: فلا ينحرف. (د) من أجل $E > cB$ يفوق الانجراف المطلوب
السرعة $c$: فلا وجود لمرجع كهذا؛ بل يوجد بدلًا منه مرجع فيه $E$ صرف
(في التمرين السابق) وتتسارع الشحنة بلا حدّ
على منحى المجال.
\end{solution}

\begin{solution}{exo:b3:covariant-electromagnetism:8}
(أ) الطور $\omega t - \vect k\cdot\vect r$ يحصي أحداث عبور
القمم، وهو عدد ثابت؛ ويساوي $k^\mu x_\mu$، وثبات
الجداء من أجل كل $x^\mu$ يفرض أن يتحوّل $k^\mu$
تحوّل الرباعيات. (ب) $\omega'/c = \gamma(\omega/c -
\beta k_x)$ مع $k_x = (\omega/c)\cos\theta$؛ ومن أجل $\theta = 0$،
$\omega' = \omega\sqrt{(1-\beta)/(1+\beta)}$. (ج) $\cos\theta' =
k'_x/(\omega'/c)$: وهي الصيغة المذكورة. (د) $\beta = \num{e-4}$:
أي $20.5''$؛ وعلى مدى سنة يكنس الموضع الظاهري قطعًا ناقصًا نصف محوره
تلك الزاوية --- وهو زيغ برادلي.
\end{solution}

\begin{solution}{exo:b3:covariant-electromagnetism:9}
(أ) في مرجع السكون هو مجال كولوم؛ والتعزيز يضرب المركبات
المستعرضة في $\gamma$ (فعند المستوي المستعرض، $E'_\perp = \gamma
q/4\pi\varepsilon_0b^2$) بينما المجال على منحى الحركة، مقدَّرًا
أمامها أو خلفها على مسافة مخبرية $r$، يقابل مسافة $\gamma r$ في مرجع السكون:
فيصغر بالعامل $1/\gamma^2$. (ب) تمرّ الفطيرة ذات العرض الزاوي
$1/\gamma$ بالسرعة $v$: فمدتها $\sim (b/\gamma)/v$. (ج)
$E' = 7250 \times \num{1.6e-19} \times \num{9e9}/\num{e-4} \approx
\qty{0.1}{V/m}$، وتدوم $2b/\gamma c \approx \qty{9}{fs}$. (د)
$B' = vE'/c^2 \approx E'/c$ وكلا الثابتين من الرتبة $O(1/\gamma^2)$:
فهي موضعيًا لا تُميَّز من نبضة ضوء --- وهو أساس
وصف ``الفوتون المكافئ'' لتصادمات الشحنات السريعة.
\end{solution}

\begin{solution}{exo:b3:covariant-electromagnetism:10}
(أ) $E'_y = \gamma(E - vB) = \gamma E_0(1 - \beta)\cos(\cdots) =
E_0\sqrt{(1-\beta)/(1+\beta)}\cos(\cdots)$، و $B'_z = E'_y/c$ بالجبر
نفسه. (ب) يتحوّل الطور بعامل دوبلر
نفسه: أي الموجة المعدومة نفسها، أحمر وأضعف. (ج) تهبط كثافة الطاقة
بمربع عامل دوبلر؛ ولا يتغير عدد الفوتونات ---
إذ يحمل $h\nu$ لكل فوتون العامل كله. (د) ``الركوب
بجانبها'' يعني أن يكون العامل $\to 0$ لمّا $\beta \to 1$: فالموجة
لا تصير ساكنة أبدًا لأن ثابتيها معدومان --- فلا
مرجع يقف فيه الضوء، وهو ما بدأت منه نسبية هذا الكتاب.
\end{solution}

\begin{solution}{exo:b3:covariant-electromagnetism:11}
(أ) في كل لفة، $\Delta E = P \times 2\pi r/c$: وهي النتيجة المذكورة. (ب)
$\gamma = \num{e11}/\num{5.11e5} = \num{1.96e5}$: ومنه $\Delta E \approx
\qty{2.9}{GeV}$ في اللفة --- أي ثلاثة بالمئة من طاقة الحزمة
تُحقن من جديد في كل دورة؛ وكانت كليسترونات ذلك المصادم أكبر جهاز
إرسال راديوي على الأرض. (ج) البروتونات: $\gamma = 7250$، و $\gamma^4$
أصغر بالعامل $(m_{\text{p}}/m_{\text{e}})^4 \approx \num{e13}$: أي نحو
\qty{5}{keV} في اللفة --- وهو مهمل. (د) $\gamma^4 = (E/mc^2)^4$:
فعند الطاقة نفسها يشعّ الإلكترون $\num{e13}$ مرة أكثر؛ ومن هنا
المصادمات الخطية (إذ تشعّ مرة لا في كل لفة) أو الحلقات العملاقة
للإلكترونات، بينما تتسع حلقات البروتونات إلى \qty{27}{km} بلا عناء.
\end{solution}

\begin{solution}{exo:b3:covariant-electromagnetism:12}
(أ) $|\vect p|$ ثابتة (فلا شغل)؛ و $\dot{\vect p} = q\vect
v\wedge\vect B$ يدير $\vect p$ بالمعدل $qvB/p = qB/\gamma m$. (ب)
الانزلاق في كل لفة $2\pi(\gamma - 1)$؛ ومع نموّ $\gamma - 1$ خطيًا
حتى $\Gamma$، يكون الانزلاق الكلي $\approx \pi N\Gamma$؛ وربع الدور
هو $\pi/2$: ومنه $\Gamma \approx 1/2N$. ومن أجل $N = 100$: $\Gamma =
\num{5e-3}$ و $E_k \approx \qty{4.7}{MeV}$ --- وهو المقياس الصحيح؛ وقد
مدّته الآلات الحقيقية إلى $\sim\qty{20}{MeV}$ بجهود عالية على الأقطاب
(أي لفات أقلّ). (ج) السنكروسيكلوترون: نخفض التواتر الراديوي في
أثناء كل نبضة ليتابع $qB/\gamma m$. والسيكلوترون المتساوي الأزمنة: ندع
$B(r)$ ينمو $\propto\gamma(r)$ فلا تتغير النسبة أبدًا وتبقى
الحزمة مستمرة. (د) $\gamma = 1 + 590/938 = 1.63$: فيجب أن يفوق مجال
الحافة المجال المركزي بذلك العامل.
\end{solution}

\begin{solution}{pb:b3:covariant-electromagnetism:1}
\textbf{1.} $u = I/nSe = \qty{7.4e-4}{m/s}$؛ و $u^2/c^2 =
\num{6e-24}$.
\textbf{2.} $\lambda = I/u = nSe = \qty{1.36e4}{C/m}$ --- أي أربعة
عشر كيلوكولومًا في المتر في كل تيار.
\textbf{3.} $B = \mu_0I/2\pi d = \qty{2.0e-4}{T}$؛ و $F = qvB =
\num{e-6} \times \num{e5} \times \num{2e-4} = \qty{2e-5}{N}$،
موجَّهة \emph{نحو} السلك (فالتياران المتوازيان يتجاذبان).
\textbf{4.} تتقاصّ الشحنتان الخطيتان تقاصًّا تامًّا: $\lambda_{\text{كلي}}
= 0$، فلا مجال في المرتبة الصفرية.
\textbf{5.} إلكترون زائد واحد في المتر: $E = 2k\lambda_1/d =
\qty{2.9e-7}{V/m}$، والقوة $\qty{2.9e-13}{N}$ --- أي أصغر $10^8$ مرة
من القوة المغناطيسية. فيمكن للسلك أن يخفي مئة مليون
إلكترون شارد في المتر قبل أن تنافس الكهرباء الساكنة المغناطيسية:
فنسبة القوة إلى $\vect B$ آمنة.
\textbf{6.} القوة $\qty{2e-5}{N}$ تعادل عشرين وزن حبة رمل: أي إنها سهلة
القياس.
\textbf{7.} الشبكة: $-v$. والإلكترونات (المنجرفة بالسرعة $u$ عكس $I$،
أي عكس التعزيز): سرعتها $(u + v)/(1 + uv/c^2)$ --- أي أسرع
من الشبكة.
\textbf{8.} $\lambda' = \gamma_v(\lambda_{\text{كلي}} - vI/c^2) =
-\gamma_v vI/c^2$.
\textbf{9.} $\lambda' = -\num{e5} \times 10/\num{9e16} =
\qty{-1.1e-11}{C/m}$: أي نحو $7 \times 10^7$ إلكترون زائد في
المتر. وتيار الإلكترونات، وهو الأسرع في هذا المرجع، هو الأكثف
--- إذ يتقلص كل تيار بعامل $\gamma$ الخاص به.
\textbf{10.} $E' = 2k|\lambda'|/d = 2 \times \num{9e9} \times
\num{1.1e-11}/0.01 = \qty{20}{V/m}$، مشيرًا نحو السلك؛ والقوة
$qE' = \qty{2e-5}{N}$، جاذبة.
\textbf{11.} وهي مطابقة للمقدار $qvB$ إلى العامل $\gamma_v = 1 +
\num{5.6e-8}$: أي بالضبط قانون تحوّل قوة مستعرضة
(أي الرباعية القوية)، $F_{\text{سكون}} = \gamma F_{\text{مخبر}}$.
\textbf{12.} ما زال السلك يحمل تيارًا (أكبر)، ومن ثمّ $\vect
B' \neq 0$ --- لكن شحنتنا ساكنة هنا، والمجال المغناطيسي
لا يمسك إلا بالشحنات المتحركة.
\textbf{13.} $|\lambda'|/\lambda = \gamma_v vu/c^2 = \num{8.2e-16}$.
\textbf{14.} إنه يُضرب في $\lambda/e \approx \num{8.5e22}$
شحنة أولية في المتر: $\num{8.2e-16} \times \num{8.5e22}
\approx \num{7e7}$ إلكترون في المتر --- وهو مقدار عياني. فالمادة
مشحونة شحنًا هائلًا إلى حدّ يجعل تعادلها حدَّ موسى؛ والنسبية
تميل بالموسى.
\textbf{15.} $F/L = \mu_0I_1I_2/2\pi d = \num{2e-7} \times 100/0.01 =
\qty{2e-3}{N/m}$ --- وهي الصيغة التي \emph{عرّفت} الأمبير يومًا.
\textbf{16.} $B = \mu_0nI = 4\pi\times10^{-7} \times \num{e4} \times
10 = \qty{0.13}{T}$؛ والضغط المغناطيسي $B^2/2\mu_0 \approx
\qty{6300}{Pa} \approx \qty{0.6}{N/cm^2}$: فسيارات بأكملها معلقة
على نسبية مُكامَلة.
\textbf{17.} $c \to \infty$ يقتل $\lambda'$، ومعه القوة،
والجزء المغناطيسي من المجال، والمحركات والمولدات ورافعات مكبّات الخردة:
فلا وجود لامغناطيسية خارج النسبية.
\textbf{18.} في السكون ترى الشحنة سلك المخبر متعادلًا: فيتقاصّ
التياران بأيّ دقة تكون عليها المادة متعادلة --- وهي معروفة
تجريبيًا بدقة مذهلة --- فلا تؤثر أيّ قوة.
\textbf{19.} أن صورة ``المنابع المعطاة'' كانت شريحة مرجع
واحد من موتّر واحد: فالمجال $\vect B$ هو الجزء من المجال
الكهرطيسي الذي تسنده حركة راصد بعينه إلى
التيارات لا إلى الشحنات.
\textbf{20.} سبين الإلكترون: فكل إلكترون مغنط أوّلي
(أي ``تيار'' ذاتي كمومي)، والحديد مادة تصطفّ فيها
هذه المغانط --- ويتناول ذلك \cref{ch:b3:spin-two-level} و
\cref{ch:b3:electromagnetism-in-matter}.
\textbf{21.} في مرجع سكون الحزمة يكون التنافر كولوميًا صرفًا
و\emph{أعظميًا}؛ وفي المخبر يكاد التجاذب المغناطيسي للتيارين
المتوازيين يلغيه. ولا تناقض: فالقوة المستعرضة في المخبر
هي قوة مرجع السكون مقسومة على $\gamma$، ويكمل الحساب
تباطؤ الديناميك (أي \omterm{prop:b3:relativistic-kinematics:dilation}{تمدد الزمن}).
\textbf{22.} القوة المحصلة في المخبر $= qE_{\text{مخبر}} - qvB_{\text{مخبر}} =
(1 - \beta^2)qE_{\text{مخبر}} = F_{\text{كولوم}}/\gamma^2$ بحسب
التحويل؛ وفي مرجع السكون يكون العامل شفافًا: فكهرباء ساكنة صرفة،
مع مشاهدة التمدد عبر زمن متمدد.
\textbf{23.} متجهة بوينتينغ: إذ تنساب الطاقة عبر
\emph{المجالات} حول النواقل بسرعة تقارب $c$، والإلكترونات
تنظّمها لا غير؛ فنحاس السلك دليل لا أنبوب.
\textbf{24.} $vu/c^2 \approx \num{8e-21}$ --- ومع ذلك رأى أورستد
بوصلته تنحرف بجانب سلك سنة 1820: فمضاعِف
$10^{22}$ شحنة كان يعمل منذ قرن قبل أن يستطيع أحد
تسميته.
\textbf{25.} سلك متعادل يحمل \qty{10}{A}، منظورًا إليه من \qty{e5}{m/s}،
يحمل $\qty{-1.1e-11}{C/m}$: أي إن فرق الحساب $10^{-16}$
الناتج عن \omterm{prop:b3:relativistic-kinematics:contraction}{تقلص الأطوال}، مضروبًا في $10^{23}$ شحنة في المتر،
\emph{هو} القوة المغناطيسية --- فالمغناطيسية كلها نسبية
مدقَّقة بخطوة حلزون.
\end{solution}
